\chapter{Arrêt optimal et contrôle impulsionnel}\label{ch:qm:optimal-stopping-and-impulse-control}

Le desk de trésorerie d'une banque porte une exposition en euros qui dérive au gré des flux de ses clients~: chaque paiement,
chaque opération de financement du négoce la déplace de quelques millions. Chaque ticket de couverture coûte un montant fixe~; porter
l'exposition coûte une charge de risque qui croît avec son carré. Couvrir après chaque flux client gaspille les frais~;
ne jamais couvrir laisse le risque croître sans limite. La réponse est une bande~: ne rien faire tant que l'exposition reste
dans $\pm b$, la ramener à zéro lorsqu'elle touche le bord. Avec les chiffres du desk, la meilleure bande est
$\pm$34.6~millions de dollars, une couverture tous les trois jours en moyenne, et la politique coûte deux fois moins cher que de couvrir
une fois par jour à la clôture~; si les frais par ticket sont divisés par deux, la bande ne se resserre que de 16~\%, car la bande
optimale croît avec la racine quatrième des frais. Décider quand agir est un problème d'arrêt~; décider quand et de combien
agir, de façon répétée, relève du \omterm{def:qm:optimal-stopping-and-impulse-control:impulse}{contrôle impulsionnel}. Ce chapitre résout les deux~: l'\omterm{def:qm:optimal-stopping-and-impulse-control:snell}{enveloppe de Snell} en temps discret,
les frontières libres et le \omterm{def:qm:optimal-stopping-and-impulse-control:fbp}{recollement C1} (smooth pasting) en temps continu, le \omterm{def:qm:optimal-stopping-and-impulse-control:impulse}{contrôle impulsionnel} et sa \omterm{def:qm:optimal-stopping-and-impulse-control:notrade}{région de non-transaction}, et
le \omterm{def:qm:optimal-stopping-and-impulse-control:singular}{contrôle singulier}, sa limite lorsque les coûts sont proportionnels.

\section{Arrêt optimal en temps discret~: l'enveloppe de Snell}

\begin{definition}[Problème d'arrêt optimal, enveloppe de Snell]\label{def:qm:optimal-stopping-and-impulse-control:snell}
Étant donné un processus de gain adapté et intégrable $(G_n)_{n\le N}$, le \emph{problème d'arrêt optimal}\index{problème d'arrêt optimal}
consiste à trouver $\sup_\tau\E[G_\tau]$ sur les \omterm{def:qm:probability-at-speed:stopping}{temps d'arrêt} $\tau \le N$. Son \emph{enveloppe de Snell}\index{enveloppe de Snell}
est définie par récurrence rétrograde par $U_N = G_N$ et $U_n = \max(G_n, \E_n[U_{n+1}])$.
\end{definition}

\begin{theorem}[L'enveloppe de Snell résout le problème d'arrêt]\label{thm:qm:optimal-stopping-and-impulse-control:snell}
$U$ est la plus petite \omterm{def:qm:probability-at-speed:martingale}{surmartingale} dominant $G$~; $U_n = \sup_{\tau \ge n}\E_n[G_\tau]$~; et $\tau^* = \min\{n : U_n = G_n\}$ est
optimal.
\end{theorem}

\begin{proof}
$U_n \ge \E_n[U_{n+1}]$ et $U_n \ge G_n$ par construction. Si $Y$ est une \omterm{def:qm:probability-at-speed:martingale}{surmartingale} dominant $G$, alors $Y_N \ge
U_N$ et, par récurrence rétrograde, $Y_n \ge \max(G_n, \E_n[Y_{n+1}]) \ge U_n$. Jusqu'à $\tau^*$, l'enveloppe égale son \omterm{def:qm:probability-at-speed:condexp}{espérance
conditionnelle}, donc $U_{n\wedge\tau^*}$ est une \omterm{def:qm:probability-at-speed:martingale}{martingale} et $\E[G_{\tau^*}] = \E[U_{\tau^*}] = U_0$~; tout autre $\tau$ donne $\E[G_\tau] \le
\E[U_\tau] \le U_0$ par le théorème d'arrêt pour la \omterm{def:qm:probability-at-speed:martingale}{surmartingale}.
\end{proof}

\begin{definition}[Région de continuation, région d'arrêt]\label{def:qm:optimal-stopping-and-impulse-control:regions}
Pour un problème markovien avec $G_n = g(n, X_n)$ et $U_n = u(n, X_n)$, la \emph{région d'arrêt}\index{région d'arrêt} est
$\{(n, x) : u = g\}$ et la \emph{région de continuation}\index{région de continuation} en est le complémentaire, où attendre est
strictement préférable.
\end{definition}

Le solveur du projet en cours, au chapitre~9, calcule l'\omterm{def:qm:optimal-stopping-and-impulse-control:snell}{enveloppe de Snell} en ajoutant le gain d'arrêt
à la récurrence rétrograde~: $V_t = \max(\text{arrêt}, \text{continuation})$. Une option américaine est le problème d'arrêt le plus connu~;
le livre~5, chapitre~6, la valorise et en nomme la frontière. Ici, les mêmes mathématiques décident
quand prendre ses profits sur une transaction.

\section{Frontières libres et inéquations variationnelles}

En temps continu, avec $X$ une diffusion de générateur $\mathcal L$ et un taux d'actualisation $r$, la valeur $V(x) =
\sup_\tau\E_x[e^{-r\tau}g(X_\tau)]$ d'un problème d'arrêt perpétuel vérifie, là où elle est régulière,
\[
\max\bigl(\mathcal LV - rV,\ g - V\bigr) = 0:
\]
dans la \omterm{def:qm:optimal-stopping-and-impulse-control:regions}{région de continuation} $V > g$ et $V$ résout l'équation $\mathcal LV = rV$~; dans la \omterm{def:qm:optimal-stopping-and-impulse-control:regions}{région d'arrêt} $V = g$ et
attendre ne rapporterait pas, $\mathcal Lg - rg \le 0$.

\begin{definition}[Inéquation variationnelle, problème à frontière libre, recollement C1]\label{def:qm:optimal-stopping-and-impulse-control:fbp}
La relation ci-dessus est une \emph{inéquation variationnelle}\index{inéquation variationnelle}. Écrit sous la forme «~résoudre $\mathcal LV = rV$ sur
une région inconnue $D$ avec $V = g$ sur son bord~», c'est un \emph{problème à frontière libre}\index{problème à frontière libre}~:
la frontière $\partial D$ fait partie de la solution. La condition supplémentaire qui la fixe, $V' = g'$ sur $\partial D$ lorsque $g$ est
régulière, est le \emph{recollement C1}\index{recollement C1} (smooth pasting, ou smooth fit).
\end{definition}

\begin{proposition}[Prendre ses profits sur une position qui dérive]\label{prop:qm:optimal-stopping-and-impulse-control:drift}
Soit $X_t = x + \mu t + \sigma W_t$ le P\&L courant d'une position, $c$ le coût de clôture, et $r > 0$ un taux d'actualisation.
La règle optimale pour $\sup_\tau\E[e^{-r\tau}(X_\tau - c)]$ clôture la position au premier instant où $X$ atteint $b^* = c + 1/\theta$, où $\theta > 0$ résout
$\tfrac12\sigma^2\theta^2 + \mu\theta - r = 0$, et $V(x) = (b^* - c)e^{\theta(x - b^*)}$ pour $x < b^*$.
\end{proposition}

\begin{proof}
Pour un seuil $b$, $\E_x[e^{-r\tau_b}] = e^{\theta(x - b)}$ pour $x < b$, puisque $e^{-rt + \theta X_t}$ est une \omterm{def:qm:probability-at-speed:martingale}{martingale} exactement lorsque
$\theta$ résout le trinôme~; la valeur de la règle de seuil $b$ vaut donc $(b - c)e^{\theta(x - b)}$, maximisée en $b$ pour $b = c + 1/\theta$.
Là, $V'(b^*) = \theta(b^* - c) = 1 = g'(b^*)$~: \omterm{def:qm:optimal-stopping-and-impulse-control:fbp}{recollement C1}. Sur $x \ge b^*$, $V = x - c$ et $\mathcal Lg - rg = \mu - r(x - c) \le 0$
car $x - c \ge 1/\theta \ge \mu/r$, si bien que l'\omterm{def:qm:optimal-stopping-and-impulse-control:fbp}{inéquation variationnelle} est satisfaite et qu'un argument de vérification, comme au chapitre~9,
s'applique.
\end{proof}

Avec $\mu = 0.5$ et $\sigma = 2$ points de base par jour, un coût de clôture de 5 pb et $r = 2\%$ par jour, $\theta = 0.0351$ et $b^* = 33.5$ pb.
La récurrence rétrograde sur une grille, en observant la position une fois par jour, s'arrête à 32.2 pb et valorise la position à 8.79
pb contre 8.80 à zéro (\cref{fig:qm:optimal-stopping-and-impulse-control:stopping})~: la grille quotidienne s'arrête un peu plus tôt,
d'environ le dépassement $0.5826\,\sigma\sqrt{\Delta t} = 1.2$ pb du chapitre~2.

\begin{omfigure}
\begin{tikzpicture}
\begin{axis}[width=11cm,height=4.8cm,
  xlabel={P\&L courant de la position $x$ (pb)},ylabel={valeur (pb)},
  tick label style={font=\small},label style={font=\small},
  xmin=-10,xmax=60,ymin=0,ymax=56,grid=major,grid style={gray!25},
  legend columns=2,legend style={font=\footnotesize,at={(0.5,-0.3)},anchor=north,draw=none,fill=none,/tikz/every even column/.append style={column sep=8pt}},
  table/col sep=comma]
\addplot[omDef,very thick] table[x=x,y=value]{figdata/methods/10-optimal-stopping-and-impulse-control/stopping.csv};
\addplot[omThm,thick,dashed] table[x=x,y=payoff]{figdata/methods/10-optimal-stopping-and-impulse-control/stopping.csv};
\addplot[gray,thin,forget plot] coordinates {(33.51,0) (33.51,56)};
\node[font=\footnotesize,anchor=north west,fill=white,inner sep=1pt] at (axis cs:34.5,52) {$b^* = 33.5$};
\legend{valeur $V(x)$,{clôturer maintenant~: $x - c$}}
\end{axis}
\end{tikzpicture}

\omcaption{Prendre ses profits sur une position qui dérive ($\mu = 0.5$, $\sigma = 2$ pb par jour, coût de clôture 5 pb, actualisation 2~\% par jour). La
valeur d'attendre (trait plein) se situe au-dessus de la valeur de clôturer maintenant (tirets) jusqu'à $b^* = c + 1/\theta$, où elle la rejoint avec la même
pente~: \omterm{def:qm:optimal-stopping-and-impulse-control:fbp}{recollement C1}. Données~: forme close, vérifiée par le solveur sur grille
du chapitre.}\label{fig:qm:optimal-stopping-and-impulse-control:stopping}
\end{omfigure}

\section{Contrôle impulsionnel et régions de non-transaction}

\begin{definition}[Contrôle impulsionnel, inéquation quasi-variationnelle]\label{def:qm:optimal-stopping-and-impulse-control:impulse}
Le \emph{contrôle impulsionnel}\index{contrôle impulsionnel} agit sur un état par des interventions discrètes~: aux \omterm{def:qm:probability-at-speed:stopping}{temps d'arrêt} $\tau_1 < \tau_2 <
\dots$, il déplace $X$ par des sauts $\xi_i$, chacun payant un coût dont la partie fixe est $K > 0$. Sa \omterm{def:qm:stochastic-control:problem}{fonction valeur} vérifie une
\emph{inéquation quasi-variationnelle}\index{inéquation quasi-variationnelle}~: $\max(\mathcal LV - rV - f,\ \mathcal MV - V) = 0$ (pour une
minimisation de coût, les signes étant inversés), où $\mathcal MV(x) = \sup_\xi\{V(x + \xi) - K - \text{coût}(\xi)\}$ est la valeur d'agir
maintenant~: l'obstacle dépend de l'inconnue $V$ elle-même.
\end{definition}

\begin{definition}[Région de non-transaction]\label{def:qm:optimal-stopping-and-impulse-control:notrade}
La \emph{région de non-transaction}\index{région de non-transaction} (no-trade region) d'une politique de négociation ou de couverture est l'ensemble des états dans lesquels elle ne
négocie pas~; pour une seule exposition, c'est un intervalle $(-b, b)$, et une \emph{politique par bande} (band policy) renvoie l'exposition vers $\pm a$ lorsque
l'exposition atteint $\pm b$.
\end{definition}

Pour le desk, supposons que l'exposition évolue comme $\sigma W_t$ entre deux couvertures, que son portage coûte $\gamma X_t^2$ par unité de temps, et
que chaque ticket de couverture coûte $K$ plus $c$ par unité négociée.

\begin{proposition}[Le coût d'une politique par bande]\label{prop:qm:optimal-stopping-and-impulse-control:band}
Le coût moyen à long terme de la politique par bande $(a, b)$ est
\[
C(a, b) = \Bigl(K + c(b - a) + \frac{\gamma(b^4 - a^4)}{6\sigma^2}\Bigr)\frac{\sigma^2}{b^2 - a^2}.
\]
Avec un coût fixe seul ($c = 0$), la réinitialisation optimale est $a = 0$, et $C(0, b) = K\sigma^2/b^2 + \gamma b^2/6$ est minimisé en
\[
b^* = \Bigl(\frac{6K\sigma^2}{\gamma}\Bigr)^{1/4}, \qquad C^* = 2\sqrt{K\sigma^2\gamma/6},
\]
où les frais et la charge de risque contribuent chacun pour moitié.
\end{proposition}

\begin{proof}
Un cycle court de $\pm a$ à la première sortie de $(-b, b)$. D'après la formule de Dynkin (chapitre~4), $m(x) = \E_x[\tau]$ et $f(x) =
\E_x\int_0^\tau X_t^2\,dt$ vérifient $\tfrac12\sigma^2m^{\prime\prime} = -1$ et $\tfrac12\sigma^2f^{\prime\prime} = -x^2$ avec des valeurs au bord nulles~: $m(x) = (b^2 -
x^2)/\sigma^2$ et $f(x) = (b^4 - x^4)/(6\sigma^2)$. Chaque cycle coûte $K + c(b - a) + \gamma f(a)$~; par le théorème de renouvellement--récompense, le
coût moyen est le coût espéré par cycle rapporté à la longueur espérée d'un cycle. Avec $c = 0$, $\partial_bC = 0$ donne $b^4 =
6K\sigma^2/\gamma$, et la substitution donne $C^*$ avec des parts égales.
\end{proof}

Avec $\sigma = 20$ (millions de dollars par racine de jour), $\gamma = \text{\$}~0.5$ par (million de dollars)$^2$ par jour et $K = \text{\$}~300$,
$b^* = 34.6$ millions et $C^* = \text{\$}~200$ par jour, avec une couverture toutes les $b^{*2}/\sigma^2 = 3$ jours en moyenne
(\cref{fig:qm:optimal-stopping-and-impulse-control:exposure,fig:qm:optimal-stopping-and-impulse-control:cost}).
Couvrir une fois par jour à la clôture coûte à la place $K + \gamma\sigma^2/2 = 400$ par jour. La courbe de coût est plate près de son minimum et
raide vers les petites bandes~: une bande de 10 millions coûte 1\,208 par jour.

\begin{omfigure}
\begin{tikzpicture}
\begin{axis}[width=11cm,height=4.8cm,
  xlabel={jours},ylabel={exposition (millions de dollars)},
  tick label style={font=\small},label style={font=\small},
  xmin=0,xmax=10,ymin=-42,ymax=42,grid=major,grid style={gray!25},
  table/col sep=comma]
\addplot[omDef,thin] table[x=day,y=x]{figdata/methods/10-optimal-stopping-and-impulse-control/exposure.csv};
\addplot[omThm,dashed,thick,domain=0:10] {34.64};
\addplot[omThm,dashed,thick,domain=0:10] {-34.64};
\end{axis}
\end{tikzpicture}

\omcaption{Dix jours simulés de l'exposition du desk sous la bande optimale $\pm 34.6$ millions (pointillés)~: les flux clients la font
évoluer librement à l'intérieur de la \omterm{def:qm:optimal-stopping-and-impulse-control:notrade}{région de non-transaction}, et chaque contact avec la bande déclenche une couverture ramenant l'exposition à zéro. Données~: le tutoriel du chapitre,
à graine fixée.}\label{fig:qm:optimal-stopping-and-impulse-control:exposure}
\end{omfigure}

\begin{omfigure}
\begin{tikzpicture}
\begin{axis}[width=11cm,height=4.8cm,
  xlabel={demi-largeur de la bande $b$ (millions de dollars)},ylabel={coût (dollars par jour)},
  tick label style={font=\small},label style={font=\small},
  xmin=8,xmax=90,ymin=0,ymax=1000,grid=major,grid style={gray!25},
  legend columns=2,legend style={font=\footnotesize,at={(0.5,-0.3)},anchor=north,draw=none,fill=none,/tikz/every even column/.append style={column sep=8pt}},
  table/col sep=comma]
\addplot[omDef,very thick] table[x=b,y=cost]{figdata/methods/10-optimal-stopping-and-impulse-control/cost.csv};
\addplot[omThm,dashed,thick] table[x=b,y=daily]{figdata/methods/10-optimal-stopping-and-impulse-control/cost.csv};
\addplot[black,only marks,mark=*,mark size=2pt] coordinates {(34.64,200)};
\legend{politique par bande $K\sigma^2/b^2 + \gamma b^2/6$,couverture quotidienne à la clôture}
\end{axis}
\end{tikzpicture}

\omcaption{Coût quotidien moyen de la politique par bande en fonction de la demi-largeur de la bande, pour les chiffres du desk ($K = 300$,
$\sigma = 20$, $\gamma = 0.5$)~: minimum de 200~\$ pour $b^* = 34.6$ millions (point), soit la moitié du coût d'une couverture quotidienne à la clôture. Données~:
forme fermée, vérifiée par simulation dans le tutoriel.}\label{fig:qm:optimal-stopping-and-impulse-control:cost}
\end{omfigure}

\section{Contrôle singulier et réflexion}

Lorsque le coût est seulement proportionnel ($K = 0$), les petites transactions coûtent peu, et la politique optimale négocie en continu mais seulement
à la frontière.

\begin{definition}[Contrôle singulier, mouvement brownien réfléchi]\label{def:qm:optimal-stopping-and-impulse-control:singular}
Un \emph{contrôle singulier}\index{contrôle singulier} agit par un processus à variation finie qui ne croît que sur un ensemble
de temps de mesure de Lebesgue nulle. Le \emph{mouvement brownien réfléchi}\index{mouvement brownien réfléchi} (reflected Brownian motion) sur $[-b, b]$ est $X =
\sigma W + L - U$, avec $L$ et $U$ non décroissants, n'augmentant que lorsque $X = -b$ et $X = b$ respectivement~: la poussée minimale qui maintient $X$
dans l'intervalle.
\end{definition}

Le \omterm{def:qm:optimal-stopping-and-impulse-control:singular}{mouvement brownien réfléchi} sur $[-b, b]$ suit la \omterm{def:qm:stochastic-differential-equations:kolmogorov}{loi stationnaire} uniforme, donc $\E[X^2] = b^2/3$, et ses bords poussent à un taux total
à long terme de $\sigma^2/(2b)$ (en appliquant la formule d'Itô à $X^2$~: la dérive $\sigma^2$ est compensée par $2b$ fois le taux de poussée). Le
coût moyen est $c\sigma^2/(2b) + \gamma b^2/3$, minimisé en
\[
b^* = \Bigl(\frac{3c\sigma^2}{4\gamma}\Bigr)^{1/3}:
\]
la loi en racine cubique pour les coûts proportionnels, contre la racine quatrième pour les coûts fixes
(\cref{fig:qm:optimal-stopping-and-impulse-control:laws}). Avec $c = \text{\$}~25$ par million (un quart de point de base),
$b^* = 24.7$ millions et le coût est de 304~\$ par jour. Avec les deux coûts, la bande optimale se déclenche à 42.4 millions et renvoie
à 12.8 millions, et non à zéro, pour 418~\$ par jour~: le coût fixe rend les transactions grandes, le coût proportionnel les fait s'arrêter plus tôt.
Davis et Norman (1990) ont résolu la version de portefeuille, avec une \omterm{def:qm:optimal-stopping-and-impulse-control:notrade}{région de non-transaction} autour de la \omterm{def:qm:stochastic-control:fraction}{fraction de Merton} du chapitre~9.

\begin{omfigure}
\begin{tikzpicture}
\begin{axis}[width=10cm,height=4.8cm,xmode=log,ymode=log,
  xlabel={coût relatif à celui du desk ($K/300$ ou $c/25$)},ylabel={bande optimale (millions de dollars)},
  tick label style={font=\small},label style={font=\small},
  xmin=0.1,xmax=10,ymin=9,ymax=70,grid=major,grid style={gray!25},
  ytick={10,20,40,60},yticklabels={10,20,40,60},
  legend columns=2,legend style={font=\footnotesize,at={(0.5,-0.3)},anchor=north,draw=none,fill=none,/tikz/every even column/.append style={column sep=8pt}},
  table/col sep=comma]
\addplot[omDef,very thick] table[x=scale,y=b_fixed]{figdata/methods/10-optimal-stopping-and-impulse-control/laws.csv};
\addplot[omThm,very thick] table[x=scale,y=b_prop]{figdata/methods/10-optimal-stopping-and-impulse-control/laws.csv};
\legend{frais fixes~: pente $1/4$,coût proportionnel~: pente $1/3$}
\end{axis}
\end{tikzpicture}

\omcaption{Politique optimale de bande selon l'ampleur du coût de transaction, en échelles logarithmiques~: avec des frais fixes par ticket, la bande croît comme la
racine quatrième des frais, avec un coût proportionnel comme la racine cubique. Une multiplication par cent du coût déplace les bandes d'un facteur
3.2 et 4.6. Données~: formes fermées, tutoriel du chapitre.}\label{fig:qm:optimal-stopping-and-impulse-control:laws}
\end{omfigure}

\section{Tutoriel~: la bande de couverture}\label{tut:qm:optimal-stopping-and-impulse-control:band}

\begin{tutorial}
\textbf{Objectif.} Calculer la bande optimale du desk, la vérifier en simulant l'exposition, trouver la bande combinée avec des coûts fixes et
proportionnels, et résoudre le problème de prise de bénéfice par \omterm{def:qm:optimal-stopping-and-impulse-control:fbp}{recollement C1} (smooth pasting) et sur une grille. \textbf{État final~:}
\cref{fig:qm:optimal-stopping-and-impulse-control:exposure,fig:qm:optimal-stopping-and-impulse-control:cost,fig:qm:optimal-stopping-and-impulse-control:laws,fig:qm:optimal-stopping-and-impulse-control:stopping}
et les nombres du problème du week-end.
\begin{tutsteps}
\item \textbf{Le coût de la bande} par renouvellement--récompense, et les deux optima en forme fermée.
\omcode{code/firm/impulse/firm_impulse.py}{26}{50}{Le coût moyen d'une politique de bande, de la réflexion, et les bandes en racine quatrième et en racine cubique.}\label{lst:qm:optimal-stopping-and-impulse-control:band}
\item \textbf{Le seuil d'arrêt} par \omterm{def:qm:optimal-stopping-and-impulse-control:fbp}{recollement C1}.
\omcode{code/firm/impulse/firm_impulse.py}{95}{100}{Le seuil de prise de bénéfice pour une position en dérive.}\label{lst:qm:optimal-stopping-and-impulse-control:stop}
\item \textbf{Exécuter} \texttt{problem()}, \texttt{dp\_stopping\_check()} et \texttt{fig\_bands.py}~; la simulation de 20\,000 jours à
390 contrôles par jour donne 197~dollars par jour pour la bande optimale, et 200 avec des contrôles quatre fois plus fréquents.
\end{tutsteps}
\textbf{Pour aller plus loin.} Faire revenir les flux clients vers la moyenne (l'exposition suit un \omterm{def:qm:stochastic-differential-equations:ou}{processus d'Ornstein--Uhlenbeck}) et trouver comment la bande
s'élargit~; couvrir deux devises corrélées par une seule bande sur leur combinaison pondérée par le risque.
\end{tutorial}

\section{Construction~: bandes de couverture}\label{bld:qm:optimal-stopping-and-impulse-control:impulse}

\begin{build}
\bfield{Rôle} Décider quand les hedgers, les rééquilibreurs et les gestionnaires d'inventaire de la société échangent~: chaque couverture récurrente passe par une
bande dont la largeur est calculée, et non estimée au jugé.
\bfield{Interface} \texttt{band\_cost(b, a, sigma, gamma, K, c)}~; \texttt{reflect\_cost(b, sigma, gamma, c)}~;
\texttt{optimal\_fixed\_band}~; \texttt{optimal\_proportional\_band}~; \texttt{optimal\_band(sigma, gamma, K, c)} renvoyant $(a^*, b^*,
C^*)$~; \texttt{simulate\_band}~; \texttt{stop\_threshold\_drift(mu, sigma, r, cost)}~; \texttt{golden}.
\bfield{Règles} Unités précisées au point d'appel (exposition, temps, devise)~; les formes fermées sont utilisées là où elles existent et vérifiées par
simulation dans les tests~; la bande combinée est recherchée à la fois sur le déclencheur et sur la remise à zéro.
\bfield{Tests d'acceptation} \texttt{code/firm/impulse/tests/}~: la formule de renouvellement contre la simulation~; les lois en racine quatrième et en racine cubique~;
l'optimum numérique se réduit aux formes fermées lorsqu'un coût s'annule~; le \omterm{def:qm:optimal-stopping-and-impulse-control:fbp}{recollement C1} tient au seuil d'arrêt.
\bfield{Pour aller plus loin} Expositions qui reviennent vers la moyenne~; plusieurs devises avec une matrice de corrélation~; la \omterm{def:qm:optimal-stopping-and-impulse-control:notrade}{région
de non-transaction} du portefeuille de Davis--Norman.
\end{build}

\begin{omsources}
\item J.~L.~Snell, ``Applications of martingale system theorems'', \emph{Transactions of the AMS} 73, 1952.
\item G.~M.~Constantinides and S.~F.~Richard, ``Existence of optimal simple policies for discounted-cost inventory and cash
management in continuous time'', \emph{Operations Research} 26, 1978.
\item M.~H.~A.~Davis and A.~R.~Norman, ``Portfolio selection with transaction costs'', \emph{Mathematics of Operations Research}
15, 1990.
\item A.~E.~Whalley and P.~Wilmott, ``An asymptotic analysis of an optimal hedging model for option pricing with transaction
costs'', \emph{Mathematical Finance} 7, 1997.
\item G.~Peskir and A.~Shiryaev, \emph{Optimal Stopping and Free-Boundary Problems}, Birkhäuser, 2006.
\end{omsources}

\section{Exercices}

\begin{exercise}[$\star$]\label{exo:qm:optimal-stopping-and-impulse-control:1}
Avec $K = 300$, $\sigma = 20$ et $\gamma = 0.5$, calculer la bande optimale, son coût quotidien et le nombre moyen de couvertures par jour.
\end{exercise}

\begin{exercise}[$\star$]\label{exo:qm:optimal-stopping-and-impulse-control:2}
Sur l'arbre à deux pas de la \cref{fig:qm:probability-at-speed:tree}, calculer l'\omterm{def:qm:optimal-stopping-and-impulse-control:snell}{enveloppe de Snell} de $G_n = (S_n - 100)^+$ et dire où il
est optimal d'arrêter.
\end{exercise}

\begin{exercise}[$\star$]\label{exo:qm:optimal-stopping-and-impulse-control:3}
«~Le jour où l'exposition atteint sa plus grande valeur absolue de ce mois~» est-il un \omterm{def:qm:probability-at-speed:stopping}{temps d'arrêt} valide pour couvrir au sens de ce chapitre~?
\end{exercise}

\begin{exercise}[$\star\star$]\label{exo:qm:optimal-stopping-and-impulse-control:4}
Si les frais par ticket doublent, de combien changent la bande optimale et son coût~? Et si la volatilité des flux double~?
\end{exercise}

\begin{exercise}[$\star\star$]\label{exo:qm:optimal-stopping-and-impulse-control:5}
Avec seulement un coût proportionnel de 25~dollars par million, calculer la bande réfléchissante et son coût.
\end{exercise}

\begin{exercise}[$\star\star$]\label{exo:qm:optimal-stopping-and-impulse-control:6}
Démontrer $b^* = c + 1/\theta$ dans la \cref{prop:qm:optimal-stopping-and-impulse-control:drift} à partir du seul \omterm{def:qm:optimal-stopping-and-impulse-control:fbp}{recollement C1}, et l'évaluer pour
$\mu = 0.5$, $\sigma = 2$, $r = 0.02$, $c = 5$.
\end{exercise}

\begin{exercise}[$\star\star\star$]\label{exo:qm:optimal-stopping-and-impulse-control:7}
\emph{Programmation.} Simuler l'exposition du desk avec la bande optimale, en contrôlant 390 puis 1\,560 fois par jour, et comparer le coût moyen à la formule.
Pourquoi la bande contrôlée à intervalles discrets est-elle moins coûteuse~?
\end{exercise}

\begin{exercise}[$\star\star\star$]\label{exo:qm:optimal-stopping-and-impulse-control:8}
\emph{Repérer l'erreur.} «~Nous couvrons dès que l'exposition dépasse 10~millions~: une bande plus serrée signifie moins de risque, et moins de risque est
toujours préférable.~»
\end{exercise}

\section{Problème~: la bande de couverture}

\begin{problem}[{Problème du week-end --- quand couvrir une exposition en dérive}]\label{pb:qm:optimal-stopping-and-impulse-control:1}
L'exposition en devise $X$ d'un desk de trésorerie évolue avec les flux clients selon $\sigma W_t$, $\sigma = 20$~millions par racine carrée de jour. La détenir coûte
$\gamma X^2$ par jour avec $\gamma = 0.5$ dollar par (million)$^2$~; un ticket de couverture coûte $K = 300$ dollars~; ensuite, la transaction coûte aussi $c = 25$ dollars
par million.

\medskip\noindent\textbf{Partie I --- Les frais fixes.}
\begin{enumerate}
\item Écrire le coût moyen de la politique de bande qui ramène à zéro en $\pm b$.
\item Quelle est la demi-largeur optimale~?
\item Combien coûte-t-elle par jour~?
\item À quelle fréquence le desk couvre-t-il~?
\item Que coûte une couverture quotidienne à la clôture, et quelle est l'économie réalisée par la bande~?
\end{enumerate}

\medskip\noindent\textbf{Partie II --- Changements d'échelle.}
\begin{enumerate}[resume]
\item Si les frais sont divisés par deux, quelle est la nouvelle bande et quel est le nouveau coût~?
\item Si la volatilité des flux double~?
\item Quelle part du coût optimal représente la charge de risque~?
\item Que donne la simulation à 390 et à 1\,560 contrôles par jour~?
\item Pourquoi la bande contrôlée à intervalles discrets sort-elle légèrement moins coûteuse~?
\end{enumerate}

\medskip\noindent\textbf{Part III --- Proportional costs.}
\begin{enumerate}[resume]
\item Avec le seul coût proportionnel, quelles sont la bande et le coût~?
\item Avec les deux coûts, quels sont le déclencheur, le retour et le coût~?
\item Pourquoi le desk ne revient-il plus à zéro lors d'une transaction~?
\item Qu'est-ce qui distingue le \omterm{def:qm:optimal-stopping-and-impulse-control:singular}{contrôle singulier} de l'item~11 du \omterm{def:qm:optimal-stopping-and-impulse-control:impulse}{contrôle impulsionnel} de l'item~12~?
\item Résoudre le problème de prise de bénéfices de la \cref{prop:qm:optimal-stopping-and-impulse-control:drift} avec les chiffres du chapitre, sous forme fermée et sur une grille quotidienne.
\end{enumerate}

\medskip\noindent\textbf{Partie IV --- Jugement.}
\begin{enumerate}[resume]
\item Un terme de charge de risque quadratique est-il un modèle raisonnable du risque du desk~?
\item Si les flux clients sont à \omterm{def:qm:stochastic-differential-equations:ou}{retour à la moyenne}, la bande doit-elle être plus large ou plus étroite~?
\item Pourquoi le desk doit-il compenser ses devises avant de les encadrer par une bande~?
\item Énoncer le \emph{résultat nommé}~: la bande optimale et sa dépendance en les frais.
\item En une phrase~: pourquoi ne rien faire à l'intérieur d'une bande est-il optimal~?
\end{enumerate}
\end{problem}

\section{Questions d'entretien}

\begin{interviewq}[$\star$ \iqroles{trader, banque}]\label{iq:qm:optimal-stopping-and-impulse-control:1}
Vous devez couvrir une exposition qui dérive avec les flux clients, et chaque couverture coûte des frais fixes. Décrire la politique optimale.
\end{interviewq}

\begin{interviewq}[$\star\star$ \iqroles{chercheur}]\label{iq:qm:optimal-stopping-and-impulse-control:2}
Qu'est-ce que le \omterm{def:qm:optimal-stopping-and-impulse-control:fbp}{recollement C1} (smooth pasting), et pourquoi la \omterm{def:qm:stochastic-control:problem}{fonction valeur} doit-elle rejoindre le gain avec la même pente~?
\end{interviewq}

\begin{interviewq}[$\star\star$ \iqroles{chercheur, trader}]\label{iq:qm:optimal-stopping-and-impulse-control:3}
Pourquoi une bande de couverture optimale varie-t-elle comme la racine quatrième des frais fixes, mais comme la racine cubique d'un coût proportionnel~?
\end{interviewq}

\begin{interviewq}[$\star\star$ \iqroles{banque, chercheur}]\label{iq:qm:optimal-stopping-and-impulse-control:4}
Qu'est-ce que l'\omterm{def:qm:optimal-stopping-and-impulse-control:snell}{enveloppe de Snell}, et quel est son lien avec le pricing d'une option américaine~?
\end{interviewq}

\begin{interviewq}[$\star\star$ \iqroles{développeur}]\label{iq:qm:optimal-stopping-and-impulse-control:5}
Implémenter un hedger à bande qui ne doit pas alterner des allers-retours lorsque l'exposition se trouve sur la bande. Que faut-il prévoir dans la conception~?
\end{interviewq}

\begin{interviewq}[$\star\star\star$ \iqroles{chercheur}]\label{iq:qm:optimal-stopping-and-impulse-control:6}
Pour un \omterm{def:qm:brownian-motion:bm}{mouvement brownien} de variance $\sigma^2$ issu de 0, calculer $\E\int_0^\tau X_t^2\,dt$, où $\tau$ est la première sortie de $(-b, b)$.
\end{interviewq}
